\chapter{Implementation}
\label{ch:implementation}

\section{Development Environment}
\label{sec:dev_environment}

The experimental implementation was developed using Python in a Google Colaboratory environment running Ubuntu Linux. Table \ref{tab:software_stack} specifies the versions of key software libraries utilized throughout this study.

\begin{table}[H]
\centering
\caption{Software Packages and Version Specifications}
\label{tab:software_stack}
\begin{tabular}{lll}
\toprule
\textbf{Software / Package} & \textbf{Version} & \textbf{Primary Function} \\
\midrule
Python & 3.10.12 & Base runtime environment \\
NumPy & 1.24.3 & Numerical array operations and bitwise audits \\
Pandas & 2.0.3 & Tabular data processing and statistical aggregation \\
Scikit-learn & 1.2.2 & Machine learning algorithms, scalers, cross-validation \\
Matplotlib & 3.7.1 & Publication-quality 2D charting and visualization \\
Seaborn & 0.12.2 & Statistical correlation heatmaps and boxplots \\
\bottomrule
\end{tabular}
\end{table}

\section{Module Implementation}
\label{sec:module_implementation}

This section presents the primary implementation logic corresponding to each stage of the experimental pipeline.

\subsection{Data Ingestion and Stratified Subsampling}
The dataset is loaded directly via \texttt{scikit-learn}. A stratified subset of 10,000 samples is extracted using \texttt{StratifiedShuffleSplit} to guarantee proportional class preservation:

\begin{lstlisting}[language=Python, caption={Stratified Subsampling Implementation}]
import numpy as np
import pandas as pd
from sklearn.datasets import fetch_covtype
from sklearn.model_selection import StratifiedShuffleSplit

# Fetch dataset
data = fetch_covtype(as_frame=True)
df = data.frame

# Perform stratified subsampling to 10,000 instances
sss = StratifiedShuffleSplit(n_splits=1, train_size=10000, random_state=42)
for sample_idx, _ in sss.split(df, df['Cover_Type']):
    sample_df = df.iloc[sample_idx].copy()

X = sample_df.drop('Cover_Type', axis=1)
y = sample_df['Cover_Type']
\end{lstlisting}

\subsection{Leakage-Free Preprocessing and Scaling}
To maintain rigorous data hygiene, continuous cartographic features ($p=10$) are scaled using \texttt{StandardScaler} fit exclusively on training data, while the 44 binary indicators remain unscaled:

\begin{lstlisting}[language=Python, caption={Leakage-Free Feature Scaling}]
from sklearn.preprocessing import StandardScaler

# Identify continuous vs binary columns
continuous_cols = [
    'Elevation', 'Aspect', 'Slope', 
    'Horizontal_Distance_To_Hydrology', 'Vertical_Distance_To_Hydrology',
    'Horizontal_Distance_To_Roadways', 'Hillshade_9am', 'Hillshade_Noon',
    'Hillshade_3pm', 'Horizontal_Distance_To_Fire_Points'
]
binary_cols = [c for c in X.columns if c not in continuous_cols]

def preprocess_splits(X_tr, X_te, cont_cols, bin_cols):
    scaler = StandardScaler()
    # Fit strictly on training partition
    X_tr_cont_scaled = scaler.fit_transform(X_tr[cont_cols])
    X_te_cont_scaled = scaler.transform(X_te[cont_cols])
    
    # Recombine scaled continuous with raw binary columns
    X_tr_final = np.hstack([X_tr_cont_scaled, X_tr[bin_cols].values])
    X_te_final = np.hstack([X_te_cont_scaled, X_te[bin_cols].values])
    
    return X_tr_final, X_te_final
\end{lstlisting}

\subsection{Multi-Seed Sensitivity Experiment Harness}
The sensitivity harness evaluates algorithm performance across five distinct train-test split seeds, tracking metric variance:

\begin{lstlisting}[language=Python, caption={Split Seed Sensitivity Harness}]
from sklearn.model_selection import train_test_split
from sklearn.metrics import accuracy_score, f1_score

split_seeds = [0, 1, 42, 99, 123]
results = []

for seed in split_seeds:
    X_tr, X_te, y_tr, y_te = train_test_split(
        X, y, test_size=0.20, random_state=seed, stratify=y
    )
    X_tr_proc, X_te_proc = preprocess_splits(
        X_tr, X_te, continuous_cols, binary_cols
    )
    
    for name, model in models.items():
        model.fit(X_tr_proc, y_tr)
        preds = model.predict(X_te_proc)
        
        results.append({
            'Model': name,
            'Split_Seed': seed,
            'Accuracy': accuracy_score(y_te, preds),
            'Macro_F1': f1_score(y_te, preds, average='macro')
        })
\end{lstlisting}

\subsection{Cross-Validation and Bitwise Reproducibility Verification}
The cross-validation harness benchmarks variance across shuffle seeds ($S \in \{11, 22, 33\}$), while the reproducibility module validates deterministic output across replicated executions:

\begin{lstlisting}[language=Python, caption={Deterministic Reproducibility Verification}]
# Run 1: Execution with fixed seed 42
rf1 = RandomForestClassifier(n_estimators=100, random_state=42)
rf1.fit(X_tr_proc, y_tr)
preds1 = rf1.predict(X_te_proc)

# Run 2: Exact replicated execution with fixed seed 42
rf2 = RandomForestClassifier(n_estimators=100, random_state=42)
rf2.fit(X_tr_proc, y_tr)
preds2 = rf2.predict(X_te_proc)

# Bitwise verification
is_identical = np.array_equal(preds1, preds2)
mismatch_count = np.sum(preds1 != preds2)
print(f"Reproducibility Verified: {is_identical} (Mismatches: {mismatch_count})")
\end{lstlisting}

\subsection{Dataset Size Scaling Pipeline Implementation}
To test model scalability, training subsets of size $2,000$, $5,000$, and $8,000$ are sampled with independent preprocessing to avoid data leakage:

\begin{lstlisting}[language=Python, caption={Dataset Size Scaling Implementation}]
size_rows = []
for size in [2000, 5000, 8000]:
    if size == len(X_train_raw):
        X_small, y_small = X_train_raw.copy(), y_train.copy()
    else:
        X_small, _, y_small, _ = train_test_split(
            X_train_raw, y_train, train_size=size,
            random_state=42, stratify=y_train
        )
    size_preprocessor = ColumnTransformer(
        transformers=[("scale", StandardScaler(), numeric_columns)],
        remainder="drop"
    )
    X_small_proc = size_preprocessor.fit_transform(X_small)
    X_test_proc = size_preprocessor.transform(X_test_raw)

    for name, clf in [("Decision Tree", DecisionTreeClassifier(random_state=42)),
                      ("Random Forest", RandomForestClassifier(n_estimators=100, random_state=42))]:
        t0 = time.time()
        clf.fit(X_small_proc, y_small)
        pred = clf.predict(X_test_proc)
        size_rows.append({
            "Training_Size": size, "Model": name,
            "Accuracy": accuracy_score(y_test, pred),
            "F1_macro": f1_score(y_test, pred, average="macro", zero_division=0),
            "Time_sec": time.time() - t0
        })
\end{lstlisting}

\subsection{Hyperparameter Evaluation Pipeline Implementation}
Systematic sweeps for Decision Tree maximum depth and Random Forest tree count:

\begin{lstlisting}[language=Python, caption={Hyperparameter Evaluation Implementation}]
parameter_rows = []
# Decision Tree max_depth sweep
for depth in [5, 10, None]:
    model = DecisionTreeClassifier(max_depth=depth, random_state=42)
    t0 = time.time()
    model.fit(X_train, y_train)
    pred = model.predict(X_test)
    parameter_rows.append({
        "Model": "Decision Tree", "Setting": f"max_depth={depth}",
        "Accuracy": accuracy_score(y_test, pred),
        "F1_macro": f1_score(y_test, pred, average="macro", zero_division=0),
        "Time_sec": time.time() - t0
    })

# Random Forest n_estimators sweep
for trees in [50, 100, 200]:
    model = RandomForestClassifier(n_estimators=trees, random_state=42)
    t0 = time.time()
    model.fit(X_train, y_train)
    pred = model.predict(X_test)
    parameter_rows.append({
        "Model": "Random Forest", "Setting": f"n_estimators={trees}",
        "Accuracy": accuracy_score(y_test, pred),
        "F1_macro": f1_score(y_test, pred, average="macro", zero_division=0),
        "Time_sec": time.time() - t0
    })
\end{lstlisting}

\subsection{Feature Group Ablation Implementation}
Ablation comparing 10 terrain features against all 54 features:

\begin{lstlisting}[language=Python, caption={Feature Group Ablation Implementation}]
feature_rows = []
for group_name, cols in [
    ("Terrain features only", terrain_features),
    ("All features", X.columns.tolist())
]:
    feat_prep = ColumnTransformer(
        transformers=[("scale", StandardScaler(), cols)],
        remainder="drop"
    )
    X_tr_f = feat_prep.fit_transform(X_train_raw[cols])
    X_te_f = feat_prep.transform(X_test_raw[cols])

    for name, clf in [("Decision Tree", DecisionTreeClassifier(random_state=42)),
                      ("Random Forest", RandomForestClassifier(n_estimators=100, random_state=42))]:
        clf.fit(X_tr_f, y_train)
        pred = clf.predict(X_te_f)
        feature_rows.append({
            "Feature_Group": group_name, "Number_of_Features": len(cols),
            "Model": name, "Accuracy": accuracy_score(y_test, pred),
            "F1_macro": f1_score(y_test, pred, average="macro", zero_division=0)
        })
\end{lstlisting}

\section{Algorithm Hyperparameter Specifications}
\label{sec:hyperparameters}

To maintain experimental control and isolate seed stochasticity, hyperparameters were kept identical across all seed iterations:
\begin{itemize}
    \item \textbf{Decision Tree:} \texttt{criterion='gini'}, \texttt{splitter='best'}, \texttt{max\_depth=None}, \texttt{min\_samples\_split=2}, \texttt{min\_samples\_leaf=1}, \texttt{random\_state=42}.
    \item \textbf{Random Forest:} \texttt{n\_estimators=100}, \texttt{criterion='gini'}, \texttt{max\_features='sqrt'}, \texttt{bootstrap=True}, \texttt{n\_jobs=-1}, \texttt{random\_state=42}.
    \item \textbf{Linear SVM:} \texttt{loss='squared\_hinge'}, \texttt{penalty='l2'}, \texttt{C=1.0}, \texttt{dual=False}, \texttt{max\_iter=3000}, \texttt{random\_state=42}.
\end{itemize}

\section{Summary}
\label{sec:impl_summary}

The implementation cleanly isolates the effect of pseudo-random seeds while enforcing strict data hygiene and hyperparameter standardization. All scripts operate deterministically with verifiable outputs. Chapter \ref{ch:results} presents the quantitative results and comparative analyses produced by this implementation.
